% HD-F1-T29. GPT-6 Astra Pro; /root/transport; 22 November 2026 UTC.
\section[HD-F1-T29: uniform exclusion of the remaining Gaussian Specht sources]{HD-F1-T29: an effective uniform character bound excludes every remaining Gaussian alternating source}
\label{sec:f1-gaussian-uniform}

The residual Gaussian alternative of T24--T28 is empty.
The proof uses an effective estimate on every permutation of every
remaining partition, and then applies the actual inertia projectors.
No asymptotic character estimate is an input.

By T28, after its specified sign-conjugate intertwiner, the original
coefficient is
\[
 \rho(g)=c(g)R_\lambda(\sigma(g)),\qquad
 \lambda=(n-k,\nu),\quad
 n\ge33,\quad 4\le k<n/8,\quad \nu\vdash k .
 \tag{F1T29.1}
\]
The original scalar character and all its correlations with
$\sigma$ are retained. Write $m=\dim R_\lambda$.
We prove the following explicit bound for every nonidentity
$s\in S_n$, with $f=f(s)$ its number of fixed letters:
\[
 \boxed{\quad
 \frac{|\chi_\lambda(s)|}{m}
       \le\frac{(n-2)f+10}{n^2}.
 \quad}
 \tag{F1T29.2}
\]

\paragraph{1. The actual alternating projectors and first-row character identity.}
Let $W_\nu$ be the rational Specht module of $S_k$, of dimension
$d_\nu$. For $0\le j\le k$ use the specified subgroup
$S_j\times S_{k-j}$, and define the actual $S_{k-j}$-module
\[
 C_j=\operatorname{im}\!\left(
       \frac1{j!}\sum_{\tau\in S_j}\operatorname{sgn}(\tau)R_\nu(\tau)
                            :W_\nu\longrightarrow W_\nu\right).
 \tag{F1T29.3}
\]
The two subgroup factors commute, so the indicated image is a
module for the second factor, with $0\le d_j:=\dim C_j\le d_\nu$.
There is no replacement of the original coefficient by $C_j$.
These modules give an exact identity of its characters:
\[
 \chi_{(N-k,\nu)}
   =\sum_{j=0}^k(-1)^j
       \chi_{\operatorname{Ind}_{S_{N-k+j}\times S_{k-j}}^{S_N}
                      (\mathbf1\boxtimes C_j)}
       \qquad(N-k\ge\nu_1).
 \tag{F1T29.4}
\]

Here is a direct derivation fixing every operation in this identity.
The Frobenius characteristic sends the character of $W_\nu$
to the Schur function $s_\nu$, induction from a product subgroup
to multiplication, and the inner product of characters to the
Hall inner product
$\langle p_\mu,p_\eta\rangle=\delta_{\mu\eta}z_\mu$.
These follow from the character formula for the specified
induced module and the conjugacy-class sizes.
Frobenius reciprocity applied to (F1T29.3) therefore gives
\[
                  \operatorname{ch}(C_j)=e_j^\perp s_\nu .
 \tag{F1T29.5}
\]
The adjoint of multiplication by $p_r$ is
$r\,\partial/\partial p_r$. Consequently the generating
operator $\sum_{j\ge0}e_j^\perp z^j$ is the algebra
automorphism
\[
          p_r\longmapsto p_r+(-1)^{r-1}z^r,
      \qquad h_a\longmapsto h_a+z h_{a-1}.
 \tag{F1T29.6}
\]
Indeed the exponential generating series for $e_j$ gives
the translation in the first formula, and applying it to
$\sum h_a t^a=\exp(\sum p_rt^r/r)$ multiplies that series
by $1+zt$.

Apply (F1T29.6) to the Jacobi--Trudi determinant for $s_\nu$.
If its original column vectors are $B_1,\ldots,B_\ell$,
they become $B_i+zB_{i-1}$. A nonzero term selecting
shifted columns must shift a prefix: shifting column $i$
without shifting column $i-1$ repeats $B_{i-1}$.
Thus the coefficient of $z^j$ is precisely the minor
formed by columns $B_0,\ldots,B_{j-1},B_{j+1},\ldots,B_\ell$.
Expansion of the first row of the Jacobi--Trudi
determinant of $(N-k,\nu)$ now gives
$s_{(N-k,\nu)}=\sum_j(-1)^jh_{N-k+j}e_j^\perp s_\nu$,
which proves (F1T29.4). In particular $C_j=0$ for $j>\ell$.
This derivation uses actual sign-isotypic projectors,
actual induction and an exact character identity.

In the induced module of (F1T29.4), cosets are the
$(k-j)$-subsets of the $N$ letters. An element fixing
such a subset acts on its fibre by the corresponding
permutation on $C_j$. Its trace is the sum of these
fibre traces. Thus, for $s\in S_n$, let $a_r(s)$ be its
number of invariant $r$-subsets. Taking absolute values
in this exact fibre formula gives
\[
             \frac{|\chi_\lambda(s)|}{m}
           \le \frac{d_\nu}{m}\sum_{r=0}^k a_r(s).
 \tag{F1T29.7}
\]
All eigenvalues on $C_j$ are roots of unity, so the
absolute value of each fibre trace is at most $d_j$.

The subsets consisting entirely of fixed letters have a
more precise contribution. When $f\ge2k$, their total
in (F1T29.4), with its signs retained, is exactly
$\dim R_{(f-k,\nu)}$. The hook formula gives, for $N\ge2k$,
\[
 D(N):=\dim R_{(N-k,\nu)}
  =\binom Nk d_\nu Q(N),\qquad
 Q(N)=\prod_{i=1}^{\nu_1}
       \frac{N-k-i+1}{N-k-i+1+\nu'_i}.
 \tag{F1T29.8}
\]
Each factor of $Q(N)$ is increasing in $N$. Also
$\prod(1-x_i)\ge1-\sum x_i$ for $0\le x_i\le1$, whence
\[
 Q(n)\ge1-\frac{k}{n-2k+1}>\frac56,\qquad
 \frac{D(f)}{D(n)}
       \le\frac{\binom fk}{\binom nk}\le(f/n)^k .
 \tag{F1T29.9}
\]
The last inequality follows term by term from
$(f-i)/(n-i)\le f/n$.
Writing $b_r(s)=a_r(s)-\binom fr$, the exact fixed-letter
contribution and the remaining fibre traces consequently give
\[
 \frac{|\chi_\lambda(s)|}{m}
 \le (f/n)^k+\frac{6}{5\binom nk}\sum_{r=0}^k b_r(s)
                       \qquad(f\ge2k).
 \tag{F1T29.10}
\]

\paragraph{2. Cumulative invariant-subset counts and their actual polynomial maps.}
Put $N=n-f$ and $M=\lfloor N/2\rfloor$. The invariant-subset
generating polynomial is
\[
             (1+x)^f\prod_{\text{nontrivial cycles }\gamma}
                                     (1+x^{|\gamma|}).
\]
For polynomials with nonnegative coefficients, compare
cumulative coefficients by dividing by $1-x$.
Multiplication by a polynomial with nonnegative coefficients
preserves this comparison.
An even cycle of length $2a$ can be replaced by $a$ factors
$1+x^2$ coefficientwise; an odd cycle of length $2a+1\ge3$
can be replaced by $(1+x^3)(1+x^2)^{a-1}$ coefficientwise.
In each case the constant and highest terms agree, and
all additional coefficients on the right are nonnegative.
Two remaining three-cycle factors can be replaced by three
two-cycle factors in the cumulative comparison, since
\[
 \frac{(1+x^2)^3-(1+x^3)^2}{1-x}
      =3x^2+x^3+4x^4+4x^5+\cdots
\]
has nonnegative coefficients. A last three-cycle factor can
be replaced by one two-cycle factor in the same comparison,
because $((1+x^2)-(1+x^3))/(1-x)=x^2$.
This last operation discards one letter from the counting
model only; it does not alter the permutation or representation.
The result is the exact upper bound
\[
 \sum_{r=0}^k a_r(s)\le S_f:=
       \sum_{j=0}^{\lfloor k/2\rfloor}
          \binom Mj\sum_{i=0}^{k-2j}\binom fi.
 \tag{F1T29.11}
\]
Subtracting the common fixed-letter polynomial proves the
same inequality for $\sum b_r(s)$ with the $j=0$ term
removed from $S_f$.

For $f\ge2k$, set $q=k/(f-k+1)<1$. Successive binomial
coefficients satisfy
$\binom f{i-1}/\binom fi\le q$ for $i\le k$. Therefore
\[
 \begin{split}
 S_f&\le\frac{\binom fk}{1-q}(1+q^2)^M
              \le\frac{\binom fk}{1-q}e^{Mq^2},\\
 \sum_{r=0}^k b_r(s)
       &\le\frac{\binom fk}{1-q}\bigl(e^{Mq^2}-1\bigr).
 \end{split}
 \tag{F1T29.12}
\]
These statements include both parities of $N$ and $k$.

\paragraph{3. The middle range of the fixed-letter count.}
Define the upper counting function
$B_f=6S_f/(5\binom nk)$.
For $t=f/n\in[1/2,3/4]$, (F1T29.12) and $k/n<1/8$ give
\[
 B_f\le
 \frac65\,\frac{t-1/8}{t-1/4}
 \left[t\exp\!\left(\frac{1-t}{16(t-1/8)^2}\right)\right]^k .
 \tag{F1T29.13}
\]
The bracket is less than one: its logarithm is at most
$-(1-t)+(4/9)(1-t)<0$.
Since $k\ge4$, division by $t$ bounds $B_f/t$ by
\[
 G(t)=\frac65\,\frac{t-1/8}{t-1/4}\,
                t^3\exp\!\left(\frac{1-t}{4(t-1/8)^2}\right).
 \tag{F1T29.14}
\]
Its logarithmic derivative is
\[
 \frac{G'(t)}{G(t)}
 =\frac{6144t^4-3584t^3-160t^2+152t+3}
        {t(4t-1)(8t-1)^3}.
 \tag{F1T29.15}
\]
The numerator, on writing $y=t-1/2\ge0$, is
$6144y^4+8704y^3+3680y^2+376y-25$.
It is strictly increasing. Thus $G$ first decreases and
then increases, with a possible missing part of either
interval; its maximum is at an endpoint. The endpoint
values obey
\[
 G(1/2)=\frac9{40}e^{8/9}<\frac{27}{40}<\frac45,\qquad
 G(3/4)=\frac{81}{128}e^{4/25}
       <\frac{2025}{2688}<\frac45 .
 \tag{F1T29.16}
\]
Here $e<3$, and $e^u<(1-u)^{-1}$ for $0<u<1$ follows
by comparing the exponential and geometric series.
Consequently
\[
                  B_f\le\frac{4f}{5n}
                       \qquad(n/2\le f\le3n/4).
 \tag{F1T29.17}
\]

\paragraph{4. Both parity endpoints and all small fixed-letter counts.}
For $f_0=0$ or $1$, put $h=\lfloor k/2\rfloor$ and
$m_0=\lfloor n/2\rfloor$. From (F1T29.11),
\[
 S_{f_0}\le2\sum_{j=0}^h\binom{m_0}j
            \le\frac73\binom{m_0}h.
 \tag{F1T29.18}
\]
Indeed $h/(m_0-h+1)<1/7$ since $k<n/8$.
For even $k=2h$, the ratios
$\binom{m_0}h/\binom n{2h}$ decrease as $h$ increases
from two through the present range: the next ratio divided
by the preceding one is at most
$(2h+1)/(n-2h-1)<1$.
For odd $k=2h+1$, increasing the lower index of the
denominator from $2h$ to $2h+1$ only decreases the ratio.
Hence
\[
 B_{f_0}\le\frac{14}{5}\,
                 \frac{\binom{m_0}2}{\binom n4}
       \le\frac{42}{5(n-1)(n-3)}
       \le\frac{10}{n^2}\qquad(n\ge33).
 \tag{F1T29.19}
\]
The last inequality is equivalent to
$8n^2-200n+150\ge0$, whose expansion at $n=33+y$ is
$8y^2+328y+2262$.

The sequence $S_f$ is convex when $f$ increases by two
with its parity fixed. To see this, fix the parity
$\delta$ of $n-f$ and write $M=(n-\delta-f)/2$.
Its second difference is the coefficient of $x^k$ in
\[
 \frac{4x^2(1+x)^f(1+x^2)^{M-2}}{1-x},
 \tag{F1T29.20}
\]
which is nonnegative. This proves the assertion by
actual polynomial subtraction.
For a given integer $f<n/2$, choose $f_0\in\{0,1\}$
of the same parity and the least integer $F\ge n/2$
of that parity. Then $F<n/2+2\le3n/4$.
Linear interpolation of the convex sequence, using
(F1T29.17) and (F1T29.19), gives
\[
 B_f\le \frac{4f}{5n}+\frac{10}{n^2}.
 \tag{F1T29.21}
\]
For $f_0=1$ the needed comparison is
$F(f-1)/(F-1)\le f$, which follows from $f\le F$.
Together with paragraph 3 this proves (F1T29.21)
for every $f\le3n/4$.

\paragraph{5. Exact cancellation at large fixed-letter count.}
For $t=f/n\ge3/4$ one has $f\ge2k$, so the exact
fixed-letter contribution in (F1T29.10) is available.
Equations (F1T29.10)--(F1T29.12) give
\[
 \frac{|\chi_\lambda(s)|}{m}
       \le t^k\left[1+\frac{6}{5(1-q)}
                                      (e^{Mq^2}-1)\right].
 \tag{F1T29.22}
\]
Here $q\le1/5$ and
$Mq^2\le(4/25)k(1-t)$. Since
$1+C(e^u-1)\le e^{Cu}$ for $u\ge0$, $C\ge1$,
the right side is at most
\[
 t^k e^{(6/25)k(1-t)}
       \le t^{19k/25}\le t^2.
 \tag{F1T29.23}
\]
The middle inequality uses $1-t\le-\log t$, and
$19k/25\ge76/25>2$ proves the last one.
For a nonidentity permutation $f\le n-2$, so
$t^2\le (n-2)f/n^2$.
For $f\le3n/4$, (F1T29.7), (F1T29.21) and
$4/5\le(n-2)/n$ give the desired bound instead.
This proves (F1T29.2) in every case.

\paragraph{6. The original scalar phases, cyclic projectors and quotient curve.}
Put
\[
       a=\frac{n-2}{n^2},\qquad b_0=\frac{10}{n^2},
       \qquad \gamma=1-an-b_0=\frac{2(n-5)}{n^2}>0.
 \tag{F1T29.24}
\]
For an actual inertia $g$ write $s=\sigma(g)$, $z=c(g)$
and $d(s)=n-\#\{\text{cycles of }s\}$.
Suppose first that $s\ne1$, of order $q$.
If $z^q=1$, the actual fixed-space projector is
$q^{-1}\sum_{j=0}^{q-1}(zR_\lambda(s))^j$.
Apply (F1T29.2) to its nonidentity permutation powers.
The cyclic orbit-count identity
$q^{-1}\sum_{j=0}^{q-1}f(s^j)=n-d(s)$ gives
\[
 \begin{split}
 \frac{\dim\ker(\rho(g)-1)}m
    &\le a(n-d(s))+b_0+\frac{\gamma}{q},\\
 \frac{\operatorname{rank}(\rho(g)-1)}m
    &\ge a\,d(s)+\gamma(1-1/q)
     \ge a\,d(s)+\gamma/2 .
 \end{split}
 \tag{F1T29.25}
\]
If $z^q\ne1$, raising a fixed vector to the $q$th power
would give $v=z^qv$, so the fixed space is zero.
The same lower bound follows because
$a(n-1)+\gamma/2=1-2/n-3/n^2<1$.
If $s=1$, the actual inertia is a nontrivial scalar
and its moving rank is $m$, again at least $m\gamma/2$.
Thus (F1T29.25)'s last bound applies to every original
nonidentity inertia entry, with its actual scalar phase.

Let $C\to C/G=\mathbf P^1$ be the original full cover,
and let $H$ be the preimage under the actual permutation
homomorphism $\sigma:G\to A_n$ or $S_n$ of a letter
stabilizer. The actual quotient curve $C/H$ has degree
$n$ over $\mathbf P^1$. At each original branch value
its local permutation index is exactly $d(\sigma(g_j))$,
including zero when that permutation is the identity.
Riemann--Hurwitz on this quotient gives
\[
                  \sum_j d(\sigma(g_j))\ge2n-2.
 \tag{F1T29.26}
\]
There are at least twelve original coefficient branch
entries. Summing the actual bounds (F1T29.25) therefore
forces
\[
 2m+6=\sum_j\operatorname{rank}(\rho(g_j)-1)
     \ge m\bigl(2a(n-1)+6\gamma\bigr)
     =m\left(2+\frac{6n-56}{n^2}\right).
 \tag{F1T29.27}
\]
This is impossible. Indeed T27 gives
$m\ge n(n-1)(n-5)/6$, and for $n\ge33$ one has
$6n-56\ge4n$, so
\[
       m\,\frac{6n-56}{n^2}
        \ge\frac23(n-1)(n-5)>6.
 \tag{F1T29.28}
\]
This contradiction excludes the entire residual
Gaussian alternative (F1T29.1).

\paragraph{Framework and source scope.}
The new algebraic maps are the specified alternating
idempotents, actual product-subgroup induction, the
Frobenius characteristic with its adjoint, and the actual
inertia power projectors. The polynomial counting
comparisons are inequalities on their exact traces.
The geometric map used in the final estimate is the
actual intermediate quotient $C/H\to C/G$.
The original curve, coefficient, primitive corner and
polarization are never replaced; their Hilbert/Chow,
Poincare-biextension, cycle-class and Adams/GRR maps are
the already proved maps of T24--T28 on those same objects.
The T25 double-base comparison and its factor-two
pairing remain valid but are not needed for this stronger
contradiction on the original base.

The classical Frobenius, Jacobi--Trudi and hook formulas
are used at their exact established scope. Their
first-row adjoint identity and every inequality needed
here were proved explicitly above.
The auxiliary finite certificate checks the counting,
parity and character identities; the exclusion itself is
uniform and follows from the displayed proof.
